\section{Lo analítico en el mundo algebraico}

Veamos las consecuencias que occuren de que las variedades algebraicas son variedades complejas.

\begin{example}
    Sobre una variedad algebraica (afín o proyectiva) las funciones regulares son holomorfas y las funciones racionales son meromorfas.
\end{example}

Podriamos pensar que una función racional sea meromorfa en un punto y al mismo tiempo sea regular en el mismo punto. Sin embargo, esto no puede pasar:
\begin{theorem}
    Sea $V\subset\pp^n$ una variedad algebraica lisa y $\phi\in K(V)$. Si $f$ es holomorfa en $p$, entonces $f$ es regular en $p$.
\end{theorem}

Nos gustaría relacionar directamente las funciones regulares y holomorfas. Sin embargo, ambas nociones están definidas en topología distintas y estructuras distintas. La estructura algebraica dada por ceros de polinomios y la topología de Zariski, mientras y la estructura holomorfa (también llamada analítica) dadas por cartas holomorfas y la topología compleja. \\
Pero sí podemos demostrar la correspondencia con las funciones meromorfas y racionales definidas en toda la variedad. 

\begin{theorem}
    Sea $V\subset\pp^n$ una variedad proyectiva lisa. Entonces 
    \begin{equation*}
        K(V)=\mathcal{K}(V)
    \end{equation*}
    Es decir, las funciones racionales y meromorfas globales coinciden.
\end{theorem}

\begin{propo}
    Toda aplicación holomorfa $f:V\to W$ entre variedades algebraicas proyectivas lisas es regular.
\end{propo}

\begin{propo}
    Dos variedades algebraicas son isomorfas si, y  solo si, son biholomorfas.
\end{propo}

\begin{theorem}
    Toda variedad compleja compacta conexa dentro de una variedad algebraica lisa es una variedad algebraica lisa.
\end{theorem}